% Standalone source; the open editor compiles this file in place.
% Research status: thesis scope and first-paper decision, 2 November 2026.
\documentclass[11pt,a4paper]{article}
\usepackage[T1]{fontenc}
\usepackage[utf8]{inputenc}
\usepackage[english]{babel}
\usepackage[margin=1in]{geometry}
\usepackage{amsmath,amssymb,amsthm}
\usepackage{enumitem}
\usepackage[expansion=false]{microtype}
\usepackage[hidelinks]{hyperref}
\setlist{itemsep=2pt,topsep=4pt}
\newcommand{\E}{\mathbb{E}}
\newcommand{\R}{\mathbb{R}}
\newcommand{\norm}[1]{\left\lVert #1\right\rVert}
\title{When Test-Time Consistency Changes Predictions\\[0.4em]
\large Backward rules, prediction error, and retention in idempotent test-time training}
\author{Amit Arad\\[0.3em]
\small Advisor: Assaf Shocher (Faculty of Data and Decision Sciences)}
\date{November 2026}

\begin{document}
\maketitle
\begin{abstract}
A model can adapt to an unlabeled input by feeding its prediction back
through a frozen reference and reducing the discrepancy between the two
outputs. When does this feedback change the update, and what information
would tell us whether the changed prediction is closer to its target?

We isolate the forward value supplied to the reference from the derivative
propagated through it. In a two-update experiment with a locally trained
tabular IT$^3$ model, refreshing that value changed the second gradient and
the final prediction in five of 20 batches. All five showed a change in
the sign of a discrepancy term. On the 40 examples together, refreshing
increased mean squared error by $6.14\%$ relative to no adaptation and
$0.40\%$ relative to reusing the initial prediction.

A compatible linear example identifies what this experiment alone cannot
resolve. Two tasks can agree on every retained source label and present
the same unlabeled inputs, yet give opposite accuracy rankings to the
same update. The thesis asks what retained source information and what
assumption connecting it to the current task would suffice to certify
a useful update while preserving earlier predictions. The fixed-reference
experiment supplies the motivating case. Moving-reference stability
results provide a separate comparison under their stated assumptions.
\end{abstract}

\section{Question, setting, and scope}
The value of a consistency loss does not specify its learning dynamics:
freezing a network's parameters and detaching its output can preserve the
same forward values while changing the update direction. A moving reference
adds feedback between the two models. We first identify when the update
changes. We then ask whether the information available to the learner
determines which prediction has lower current error.

The learner starts from supervised pretraining and observes an unlabeled
stream $x_t\sim\rho_t$. The adaptive state $\theta_t$ and reference state
$\phi_t$ may persist across batches. Labels are available only to the
evaluator after pretraining. The primary quantity is current task risk
$R_t(\theta)=\E_{(x,y)\sim\mu_t}\ell(q_\theta(x),y)$, evaluated before each
batch's adaptation; post-update risk is reported separately. The initial geometry test uses the source distribution. The first dynamics
test also keeps a fixed empirical distribution stationary; covariate changes with one
unchanged labeling function and compatible recurring domains follow only
if the mechanism is supported. Arbitrary changes in hidden
labels cannot be identified from unlabeled inputs alone.

The local anonymous position manuscript \emph{Forgetting is not Catastrophic,
Our Benchmarks are} motivates evaluating the presently relevant task.
It remains private reading context; its authors and public publication
status are unverified. Its discussion of streams with labels does not
provide an unlabeled-TTA guarantee. Accordingly, old-task retention is an
objective only when those tasks remain relevant. Accuracy loss, recovery
after new adaptation, and preservation of information are different claims.

\section{What the literature and implementation already establish}
IT$^3$ introduces a dual-input predictor and already includes online
adaptation with an EMA reference; its online discussion allows deterioration
on the original distribution~\cite{it3}. Online carry-over and gradient
alignment were also studied in original TTT~\cite{ttt}. Neither online
adaptation nor the observation that aligned gradients improve task loss
is proposed as a new contribution.

The closest continual-adaptation baselines constrain the claim. CoTTA
combines weight/augmentation averaging with stochastic restoration to
pretrained weights~\cite{cotta}; EATA uses sample selection and
importance-weighted regularization against forgetting~\cite{eata}.
RMT studies consistency-loss gradient properties and teacher-based
continual adaptation~\cite{rmt}. PeTTA already considers recurring
environments and analyzes a tractable collapse model~\cite{petta}.
Recent work explicitly compares frozen and moving teachers~\cite{fixedteacher}
and combines dynamic EMA with source anchoring~\cite{desa}. Thus ``EMA can
fail,'' ``a fixed teacher helps,'' and an $A\to B\to A$ experiment are
insufficient paper claims. Nonstationary adaptation guarantees and recovery
complexity also have existing treatments~\cite{realign,recovery}.

\paragraph{What the closest analyses do and do not cover.}
The bounded audit uses PeTTA v4 (Theorem~1 and Appendix~B.3), RMT v2
(Section~3.1, Eqs.~(2)--(6)), and the fixed-teacher preprint v1
(Section~5.2, Eqs.~(3)--(5)). PeTTA bounds a scalar class-mean recurrence;
RMT compares loss derivatives with a fixed teacher probability;
the fixed-teacher paper explicitly describes error feedback through EMA.
These support rejecting generic feedback failure as a distinct contribution,
but do not directly certify our dual-input neural dynamics.
The inspected official RMT and fixed-teacher consistency paths feed images
to the teacher, without a student-prediction input. Exact revisions and code
locations are retained in the research notes; no benchmark was executed.
This limited comparison does not certify novelty~\cite{petta,rmt,fixedteacher}.

\paragraph{Version and graph audit.}
The bibliography uses the ICML 2025 IT$^3$ record. This research pass compared
arXiv v1 and v2; it did not establish byte identity between the supplied
local PDF, arXiv v2, and the proceedings PDF. The official repository was
inspected at commit \texttt{100524001d89}\footnote{Full revision:
\href{https://github.com/nikitadurasov/ittt/tree/100524001d890efc08911f3cba8f45f51f4aa7d1}{\texttt{100524001d890efc08911f3cba8f45f51f4aa7d1}}.}.
The tabular, airfoil, and age regression notebooks pass an undetached student
prediction into a reference excluded from the student optimizer. They use
absolute discrepancies. The CIFAR notebook instead reuses a shared
trainable model across the two adaptation passes and uses Jensen--Shannon
divergence. The linked \texttt{torch-ttt} implementation also reuses its
model. There is therefore no universal ``IT$^3$ backward graph.'' This
proposal selects the regression-style reference-input path and a squared
surrogate explicitly. The inspected repository did not reveal a reproducible
online EMA schedule; our ordering and optimizer-state policy must be declared
as experimental choices. A focused recheck of the pinned tabular notebook
confirms a mean over all absolute residual entries (cell 13, line 38),
an undetached reference input, and student-only Adam. Its stored model
definition uses ReLU layers, one output and default auxiliary input 50
(cells 5 and 12). Our two-output tanh model with zero auxiliary input
remains an explicit surrogate. These are static audits, not reproduced results.

\section{The specified dynamics}
Write $q_\theta(x)=f_\theta(x,0)$ and $T_\phi(x,q)=f_\phi(x,q)$. With
$r=T_\phi(x,q_\theta(x))-q_\theta(x)$, define
\begin{equation}
 D(\theta,\phi;\rho)=\tfrac12\E_\rho\norm{r}^2,
 \qquad J_q=\partial_\theta q_\theta(x).
 \label{eq:discrepancy}
\end{equation}
Holding reference parameters fixed during the student derivative retains
the reference-input Jacobian. Detaching the entire reference output removes it:
\begin{align}
 G_{\rm input}&=\E_\rho[J_q^\top(\partial_qT_\phi-I)^\top r],
 \label{eq:input}\\
 G_{\rm stop}&=-\E_\rho[J_q^\top r].
 \label{eq:stop}
\end{align}
These are chain-rule calculations for \eqref{eq:discrepancy}, not results
for every IT$^3$ implementation. $G_{\rm stop}$ is a controlled comparison.
Frozen parameters during differentiation do not imply a reference that
stays fixed over the stream.

We first study $\dot\theta=-G(\theta,\phi;\rho_t)$ and
$\dot\phi=(\theta-\phi)/\tau$, with $\tau>0$. Even for $G_{\rm input}$,
a fixed-distribution calculation gives
\begin{equation}
 \dot D=-\norm{\nabla_\theta D}^2
 +\tau^{-1}\langle\nabla_\phi D,\theta-\phi\rangle.
 \label{eq:energy}
\end{equation}
The second term has no prescribed sign. Distribution drift adds a time
derivative, or a jump in the evaluated objective at a switch. For
$G_{\rm stop}$ the first term is instead
$-\langle\nabla_\theta D,G_{\rm stop}\rangle$.

\paragraph{Discrete algorithm.}
The first study uses gradient descent and a student-first EMA,
\begin{equation}
 \theta_{n+1}=\theta_n-\eta G(\theta_n,\phi_n;\rho_n),\qquad
 \phi_{n+1}=\alpha\phi_n+(1-\alpha)\theta_{n+1},\quad 0<\alpha<1.
 \label{eq:discrete}
\end{equation}
Both states start from the pretrained model. Batch statistics are specified
and held fixed in the first comparison; there is no replay or source-label
access during adaptation. Adam, normalization updates, and alternative EMA
orderings are later robustness checks. Resetting weights without resetting
optimizer state is not treated as an equivalent baseline.

\section{Thesis scope and first-paper decision}
The thesis studies two questions in sequence. The first has a checked
answer for a fixed graph: when does a refreshed prediction change the
second update? The open question concerns a pair of compatible tasks:
what can a learner retain from the first task that certifies an improvement
on the second while preserving the first task's predictions?

\paragraph{What the experiment adds.}
IT$^3$ specifies a frozen outer model in Eq.~(6) and a moving reference in
Eq.~(7). Its discussion of Eqs.~(17)--(18) already identifies the risk of
reinforcing a wrong prediction~\cite{it3}. The audited tabular code retains
the derivative through the frozen model's prediction input. RMT compares
cross-entropy derivatives with the teacher probability held fixed
(Section~3.1, Eqs.~(2)--(4))~\cite{rmt}. Its teacher receives an image.
Our matched comparison holds optimizer history and the input derivative
rule fixed while changing the forward input. It locates the sign and
reference-derivative terms responsible for the gradient difference in
this run. That attribution is useful thesis evidence. The calculation
follows from the chain rule, and the registered accuracy test failed.
We therefore close this candidate as a stand-alone first paper.

\paragraph{Why the sign event cannot settle accuracy.}
Let $p$ be a vector of baseline predictions on $n$ examples, $d$ the
change after adaptation, and $y$ their targets. Expanding the two squared
errors gives the exact finite-step difference
\[
 R(p+d)-R(p)=\bigl(\|d\|^2-2\langle y-p,d\rangle\bigr)/n .
\]
The measured update fixes $d$. Strict improvement additionally requires
$\langle y-p,d\rangle>\|d\|^2/2$. The inequality applies to the final
predictions even when Adam transforms the gradient nonlinearly.

For a concrete limit, let $S$ contain the old feature rows, let $X$ contain
the current inputs, and suppose the learner has old labels $Sw$ and a
linear predictor $w$. Consider an update $h$ with $Sh=0$ and $d=Xh\ne0$.
Two possible fixed labeling functions have coefficients $w+h$ and $w-h$.
Each agrees with every old label. The learner sees identical old data and
unlabeled current inputs in the two worlds. With $p=Xw$, the initial
current error in both worlds is $D=\|d\|^2/n$. The update changes it to
zero in the first world and $4D$ in the second. Old predictions remain
exact in both. This construction applies to a given displacement under
the stated information budget. It makes no assumption that a trained
IT$^3$ model produces a particular linear displacement.

\paragraph{Next theorem and decision rule.}
Start with known common features, noiseless old label constraints, and an
explicit set of admissible coefficients justified before current labels
are seen. Can retained source information yield a lower bound on
$\langle y-p,d\rangle$ that exceeds the finite-step threshold? Compare
with keeping the source predictor, restoring its weights, and projecting
an update to preserve old outputs. Count the retained coefficients,
labels or summary statistics for each rule. First derive a sharp bound
in the linear model and exhibit a case in which it certifies a nonzero
update. A set that admits the baseline as a possible truth already
precludes a strict uniform gain. A result that needs hidden current
labels or reduces to an existing projection guarantee ends this paper
direction. Independent validation on a fresh checkpoint and untouched
tasks would precede any claim about trained models. Drift and recovery
require separate experiments after that decision.

\subsection{A constructed useful equilibrium and a stability distinction}
The earlier absolute model $f_W(x,z)=Wx+A_xz$ has a limitation: strict
stability of its homogeneous dynamics implies zero predictions. To test
whether a useful equilibrium is possible, consider instead the smooth
shared architecture
\begin{equation}
 f_W(x,z)=h(x)+Wx+
 A_x(z-h(x))\frac{h(x)^\top z}{\norm{h(x)}^2},
 \qquad h(x)\ne0.
 \label{eq:realization}
\end{equation}
Here $h$ is a fixed pretrained baseline; $W$ is adapted and $V$ is its EMA.
Then $q_W(x)=f_W(x,0)=h(x)+Wx$, while $f_0(x,h(x))=h(x)$ and
$\partial_z f_0(x,h(x))=A_x$. On two equally likely inputs
$x_1=(1,0)^\top$, $x_2=(3/5,4/5)^\top$, set
$h(x_1)=(1,0)^\top$ and $h(x_2)=(0,1)^\top$, with evaluation labels
equal to $h$. Thus $W=V=0$ represents correct, distinct nonzero
predictions, not collapse. Correctness is planted in the baseline;
the learner receives no labels during adaptation.

Writing $B_x=I-A_x$, the residual is
$Vx-B_xWx+O(\norm{W}^2)$. For column vectorizations $w,v$, the local
Jacobian is
\begin{equation}
 M=\begin{pmatrix}-H&K\\I/\tau&-I/\tau\end{pmatrix},
 \label{eq:block}
\end{equation}
where the two backward rules yield
\begin{align}
 H_{\rm input}&=\E[xx^\top\otimes B_x^\top B_x],&
 K_{\rm input}&=\E[xx^\top\otimes B_x^\top],\nonumber\\
 H_{\rm stop}&=\E[xx^\top\otimes B_x],&
 K_{\rm stop}&=\E[xx^\top\otimes I].
 \label{eq:blocks}
\end{align}
The covariance eigenvalues are $1/5$ and $4/5$, so current risk
$R(W)=\tfrac12\E\norm{Wx}^2$ is positive definite in $W$.
A Hurwitz $M$ therefore gives local exponential recovery to useful
predictions in this construction. Monotone risk is not implied.

\paragraph{An exact boundary and a noncommuting witness.}
Let
\[
 J=\begin{pmatrix}0&-1\\1&0\end{pmatrix},\qquad B_0=\tfrac12 I+J.
\]
For $B_{x_1}=B_{x_2}=B_0$, an exact mode calculation gives local
exponential stability of the input-gradient rule for every finite
$\tau>0$, whereas the stopped-output rule is locally exponentially
stable exactly when $\tau>10$. At $\tau=10$ its linearization has
imaginary roots. This is a linear stability crossing, not a proved
nonlinear Hopf bifurcation. The two rules are consequently more than
a scalar change of speed; noncommutation across inputs is not necessary
for this distinction.

For $B_{x_1}=B_0$ and $B_{x_2}=B_0+\tfrac1{10}\operatorname{diag}(1,-1)$,
the matrices do not commute. Exact rational Routh certificates give the
classifications below. For the student-first discrete update, use
$\eta=1/10$ and $\alpha=\tau/(\tau+\eta)$, with Jacobian
\begin{equation}
 J_{\eta,\alpha}=\begin{pmatrix}
 I-\eta H&\eta K\\
 (1-\alpha)(I-\eta H)&\alpha I+(1-\alpha)\eta K
 \end{pmatrix}.
 \label{eq:discreteblock}
\end{equation}
\begin{center}
\begin{tabular}{lrrr}
Rule & $\tau$ & $\max\Re\lambda(M)$ & $\rho(J_{\eta,\alpha})$\\ \hline
Input gradient & 1 & $-0.09933$ & $0.99024$\\
Stopped output & 1 & $+0.16592$ & $1.01767$\\
Input gradient & 20 & $-0.02134$ & $0.99786$\\
Stopped output & 20 & $-0.02156$ & $0.99782$
\end{tabular}
\end{center}
The displayed eigenvalues are rounded diagnostics. The underlying
continuous classifications use exact characteristic polynomials and
Routh signs; the discrete classifications use the exact bilinear
transformation of the characteristic polynomial and the same test.
The unstable cases have four unstable modes. No interval-wide threshold
is claimed for this noncommuting perturbation. Nonlinear simulations
from matched small student/reference perturbations agree with the
predicted recovery or growth, including a halved-step check.

\paragraph{What the construction establishes.}
A task-correct equilibrium and a directional stability distinction are
mathematically realizable in one smooth dual-input model. The result
does not show that neural pretraining learns this model or that it adapts
successfully to an unseen distribution shift. At exact initialization
it stays fixed; the experiments test recovery from parameter perturbations.
The chosen $A_x$ are not idempotent projections. A genuinely idempotent
differentiable reference on a neighborhood would satisfy $A_x^2=A_x$
at its fixed point, so measured projection-like Jacobians could invalidate
this particular mechanism. Classical feedback stability tests are not
the proposed novelty. Full derivations, exact certificates, and verification
scripts are recorded in the research notes.

\subsection{Assumptions, tracking, and rejection criteria}
A theorem beyond the constructed example must connect a smooth trained
model to a justified task-correct reference point or equilibrium family,
and control the identifiable directions. Any extension to drift must specify how an
equilibrium $z_t^\star$ moves. A bound of the form
$\norm{e(t)}\le C e^{-\gamma t}\norm{e(0)}+
C\sup_s\norm{\dot z_s^\star}/\gamma$
requires a uniform stability estimate and a justified forcing bound;
it is a target template, not a proved result here. Pointwise stable
Jacobians alone do not justify arbitrary switching. Converting state
tracking to predictive risk additionally needs local output regularity
and a task-alignment or correctness assumption.

Reject or narrow the leading direction if the stable state is useless;
if the two rules reduce to time rescaling; if the condition assumes
contraction without explaining its origin; if a correction duplicates
the closest prior work; or if predicted behavior fails inside its stated
regime in small nonlinear tests. A counterexample to a proposed certificate
is a useful result, but should not be relabeled a positive theorem.

\section{Minimal experiment and near-term milestones}
\paragraph{Completed neural geometry gate.}
A fixed five-seed experiment trained a $4\to16\to2$ tanh regressor on
256 source inputs using equal-weight zero/label input supervision inspired
by IT$^3$~\cite{it3}. It used a smooth, nonconstant two-output labeling
function, squared losses, 6000 Adam steps per seed, and 256 independent
evaluation inputs. No architecture search or checkpoint selection was
performed. The fit gate required both auxiliary-input versions to have
normalized mean squared error at most $0.05$, where normalization divides
by total evaluation-label variance.
\begin{center}
\begin{tabular}{rrrr}
Seed & Zero-input NMSE & Label-input NMSE & Fit gate\\ \hline
0 & $0.01203$ & $0.00405$ & pass\\
1 & $0.06620$ & $0.01435$ & fail\\
2 & $0.03434$ & $0.00698$ & pass\\
3 & $0.01708$ & $0.00612$ & pass\\
4 & $0.00719$ & $0.00146$ & pass
\end{tabular}
\end{center}
Four seeds passed, but no evaluation point in any seed met the prescribed
strong-rotation screen, at either predicted or true-label auxiliary input.
The screen required eigenvalues $a\pm i\omega$ with $0<a<1$ and
$\omega^2>a(1-a)$, with a numerical margin; for the normal family this is
the fast-reference separation underlying the construction. For general
nonnormal matrices it is only a screen, not a stability criterion.
The test therefore fails to support that particular explanation in this
model. It does not rule out rotation in other IT$^3$ settings.
Nor does it justify assuming projection geometry: across the four passing
seeds, median relative defects
$\norm{A_x^2-A_x}_F/(\norm{A_x}_F+10^{-12})$ range from $0.46$ to $0.94$.
All sampled prediction-input operator norms are below $0.902$, which does
not imply stability of the coupled parameter updates.

The actual full parameter fields were also differentiated, retaining
nonzero-residual terms. Both rules have positive instantaneous spectral
abscissae at the tested $\tau=1,20$, but the checkpoints have nonzero
velocity: these spectra are not equilibrium stability certificates or
predictions of sustained risk growth. Moreover, 32 of 114 student
parameters multiply the always-zero auxiliary input and receive exactly
zero gradient. With matched initialization those coordinates remain fixed
in both copies; removing them preserves the measured maximal real
eigenvalues. A strict full-parameter Hurwitz target is inappropriate.
Finite-difference gradient checks agree within $2.5\times10^{-10}$, and
halving the full-field differentiation step changes entries by less than
$1.3\times10^{-10}$. Protocol, all checkpoints, derivatives, and results
are retained in the research notes. This is a synthetic mechanism test,
not an official-implementation or distribution-shift reproduction.

\paragraph{Completed finite-time response test.}
The same five checkpoints were tested without retraining; seeds 0,2,3,4
were primary and seed 1 a declared sensitivity case. Adaptation repeatedly
used 64 unlabeled inputs, with a separate 192-input evaluator. This is a
stationary empirical-distribution stress test, not a fresh-sample stream.
Both rules were compared with fixed reference, EMA at $\tau=1,20$, and
no adaptation. The specified student-first update used $\eta=.02$ for
1000 steps ($T=20$), with one halved-step comparison at the same horizon.
Labels entered only the evaluator.

Two student-only initial perturbation directions were fixed per seed:
the student component of a local growing mode, shared across configurations,
and one seeded random direction. Both lie in the 82 non-fixed coordinates;
the reference initially remains unchanged. Prediction sensitivity was
normalized by its initial value for the same direction. The modal probe
deliberately stresses the dynamics; it is not typical noise or a worst-case
singular-value calculation. At $\tau=1$, the primary modal results are:
\begin{center}
\begin{tabular}{rrrrr}
 & \multicolumn{2}{c}{Prediction sensitivity gain}
 & \multicolumn{2}{c}{Final NMSE}\\
Seed & Input gradient & Stopped output & Input gradient & Stopped output\\ \hline
0 & $4.832$ & $5.316$ & $1.703$ & $3.810$\\
2 & $645.165$ & $51.875$ & $54.111$ & $249.731$\\
3 & $23.352$ & $63.814$ & $3.091$ & $9.749$\\
4 & $1.859$ & $1.912$ & $0.516$ & $0.672$
\end{tabular}
\end{center}
Initial primary NMSE ranges from $.00752$ to $.03673$; fixed-reference
final NMSE is $.00685$--$.03697$. The evaluator normalization differs
from the earlier 256-point geometry test. The predefined amplification
gate passes in three of four primary seeds: stopped-output fast-EMA gain
at least 2, with fixed-reference gain at most 1.25. The prediction that
stopped-output gain exceeds input-gradient gain by at least 25\% in three
seeds fails; only seed 3 meets it. Seed 2 illustrates why sensitivity and
risk rankings must remain separate. Slow EMA is less damaging over this
horizon but is not a guarantee: stopped-output seed 2 still deteriorates
materially. Fast-EMA consistency discrepancy increases in these runs;
monotone consistency improvement is not claimed.

The actual time-dependent tangent recurrence predicts paired nonlinear
prediction differences with maximum relative error $1.28\times10^{-5}$
at perturbation size $10^{-5}$. Halving the step changes final modal gains
by at most $2.40\%$ and NMSE by $0.52\%$. These checks support a finite-time
observation, not an equilibrium or ODE-limit theorem. The tangent product
is an established chain-rule tool; its successful implementation is not
novelty. All 30 configurations, including the failed-fit seed, are retained
in the research notes. No distribution drift or retention test was run.

\paragraph{Completed local feedback decomposition.}
At the same five initial checkpoints, write
$J=\partial_\theta f_\theta(x,0)$,
$L=\partial_\phi f_\phi(x,q)$ with $q$ held fixed, and
$B=I-\partial_qT_\phi$. The input-gradient zero-residual blocks of
the student velocity are $-\E[J^\top B^\top BJ]$ and
$\E[J^\top B^\top L]$; the stopped-output blocks are
$-\E[J^\top BJ]$ and $\E[J^\top L]$.
The full derivatives additionally differentiate the prefactor multiplying
the residual. These product-rule remainders were checked independently
against the saved full derivatives, with maximum entry error
$1.58\times10^{-10}$. Dropping them does not make the actual checkpoint
an equilibrium.

The frozen protocol compared full, zero-residual, and additionally
$L=J$ matrices on the 82 non-fixed coordinates per copy, for both rules
and $\tau=1,20$. A positive observable mode requires
$\Re\lambda>10^{-6}$ and
$\norm{J_{\rm eval}u}/(\norm{u}\norm{J_{\rm eval}}_2)\ge10^{-3}$.
Here $u$ is the student component of the mode.
All four primary seeds retain such a mode in every setting under both
diagnostics. Removing residual terms changes the leading observable rate
by at most $7.34\%$, passing the predefined $10\%$ comparison.
Additionally setting $L=J$ changes it by up to $34.12\%$, failing that
quantitative comparison despite preserving positive modes.

Thus neither residual terms nor derivative mismatch are necessary for
this local signature. The $L=J$ replacement is a matrix diagnostic, not
a demonstrated neural intervention; it does not explain the complete
nonlinear risk trajectory or show that mismatch is irrelevant. These
are nonstationary checkpoints, so equilibrium stability is still not
established. All 60 matrices, including the failed-fit seed, and the
independent checks are retained in the research notes.

\paragraph{Completed control: hold the prediction input fixed.}
The reference target $T_\phi(x,q_\theta(x))$ can change with the student
even when $\phi$ is fixed. Detaching its output still recomputes that
target after every update. To remove this forward path, the control
caches $c(x)=q_{\theta_0}(x)$ once and uses
\[
 r_c=T_\phi(x,c(x))-q_\theta(x),\qquad
 \dot\theta=\E[J^\top r_c].
\]
Both backward rules coincide for the control. Its initial target and
velocity equal those of the original stopped-output rule. With
$S=\partial_\theta\dot\theta$, $T=\partial_\phi\dot\theta$ and
$A=\partial_qT_\phi$, the exact initial full-block differences are
$S_{\rm stop}-S_c=\E[J^\top A J]$ and $T_{\rm stop}-T_c=0$.
The residual-weighted prefactor terms cancel at initialization; the
remaining difference has no prescribed sign. Independent differentiation
verified these identities within $6.51\times10^{-11}$.

The unchanged protocol reused all five checkpoints and the 64/192
adaptation/evaluation split. It compared original stopped-output,
original full input-gradient, and cached-input dynamics, each with fixed
reference or $\tau=1$ EMA, at $\eta=.02$ and $.01$, both to $T=20$.
All 60 trajectories completed without retraining or tuning. Labels were
evaluation-only. Every condition allocated the same 128-number prediction
cache (1024 bytes); only the control read it during updates. This is
retained information associated with the repeated inputs, not a proposed
fresh-input online method.

Let $R_0,R_e,R_c$ denote initial, original stopped-output/EMA, and
cached-input/EMA terminal NMSE. The predeclared gate required
$R_e\ge2R_0$, $R_e-R_0\ge.01$, and
$\max(R_c-R_0,0)\le.25(R_e-R_0)$ in at least three of four primary seeds
at both step sizes. All four pass. Main-step results are:
\begin{center}
\begin{tabular}{rrrrr}
Seed & Initial $R_0$ & Original $R_e$ & Cached $R_c$ & Excess removed\\ \hline
0 & $.01234$ & $3.80976$ & $.19615$ & $95.16\%$\\
2 & $.03673$ & $249.73110$ & $.11633$ & $99.97\%$\\
3 & $.01722$ & $9.74948$ & $.07419$ & $99.41\%$\\
4 & $.00752$ & $.67178$ & $.08344$ & $88.57\%$
\end{tabular}
\end{center}
The minimum reduction across both step sizes is $88.55\%$. All original
endpoints exactly reproduce the earlier baseline. Halving the step changes
terminal NMSE by at most $.517\%$ across every condition, below the $5\%$
numerical gate. The failed-fit seed 1 is retained as a sensitivity case
and also meets the numerical/effect inequalities, without entering the
primary count. All states, targets, metrics and derivative checks are
preserved in the research notes.

This supports a large effect of the forward input path under the declared
intervention. It does not establish a safe correction: every cached/EMA
primary run still ends above its initial risk. Cached/EMA discrepancy is
lower at the endpoint despite this accuracy loss; monotonicity is not
claimed. Reference targets remain exactly constant in the fixed-weight,
cached-input control; caching alone leaves weight-driven target changes.
The intervention preserves the reference update rule, not the subsequent
realized reference weights between arms. Neither a general stability
theorem nor novelty follows from the reduction.

\paragraph{Completed early prediction and its limits.}
An analysis fixed before computing the approximation tested whether
initial full derivatives predict the output difference between original
stopped-output/EMA and cached-input/EMA. This reuses the existing
trajectories after their terminal results were known; it is not an
independent experiment. Let $u,v$ denote student/reference displacements
on the 82 non-fixed coordinates, $F_0$ their common initial student
velocity, and $S,T$ each arm's full initial student/reference derivative
blocks. Starting from $u_0=v_0=0$, the forecast uses
\begin{align*}
 u_{n+1}&=u_n+\eta(F_0+Su_n+Tv_n),\\
 v_{n+1}&=\alpha v_n+(1-\alpha)u_{n+1},\qquad
 \alpha=(1+\eta)^{-1}.
\end{align*}
The predicted output separation is
$\widehat{\Delta q}_n=J_{\rm eval,0}(u_{{\rm stop},n}-u_{c,n})$.
All derivatives remain fixed at initialization; no trajectory fitting
or labels enter this calculation. The nonzero forcing $F_0$ matters:
these checkpoints are not equilibria. The first affine steps agree,
and the second student steps differ exactly by
$\eta^2(S_{\rm stop}-S_c)F_0$. This is an affine-map identity, not an
exact finite-step identity for the nonlinear model.

The fixed gate required at most $25\%$ relative RMS prediction error
at $T=2$, with non-negligible true separation, in at least three of four
primary seeds at both step sizes. All four pass:
\begin{center}
\begin{tabular}{rrrr}
Seed & Actual separation RMS, $\eta=.02$ & Error, $\eta=.02$ & Error, $\eta=.01$\\ \hline
0 & $.001403$ & $.5850\%$ & $.5836\%$\\
2 & $.005074$ & $.4531\%$ & $.4517\%$\\
3 & $.001313$ & $1.5078\%$ & $1.5039\%$\\
4 & $.001547$ & $.6380\%$ & $.6356\%$
\end{tabular}
\end{center}
All numerical checks pass: evaluator-Jacobian finite differences agree
within $4.54\times10^{-11}$, and two independent affine propagations
agree within $7.86\times10^{-14}$ relative state error. An independent
audit checks all 510 saved output comparisons. No nonlinear adaptation
or training was rerun.

This early forecast cannot be extrapolated uniformly. In primary seed 2,
the first saved error above $25\%$ occurs at $T=10.4$ at both step sizes;
by $T=20$ its relative error exceeds $240$ (over $24{,}000\%$). The
failed-fit sensitivity seed crosses the threshold at $T=8$ and is retained.
These are sampled diagnostics at intervals of $.4$, not certified exit
times. The other primary seeds remain below $25\%$ on the saved grid.
No current-risk, asymptotic-stability, retention or recovery guarantee
follows. A full-matrix affine forecast is useful evidence, but classical
linearization itself is not a novel mechanism theorem.

\paragraph{Completed absolute-loss audit: a limit on transfer.}
The local explanation changes when the squared surrogate is replaced by
the absolute discrepancy used in the selected regression path. For $N$
inputs and $d$ output coordinates, set
\[
 D_1=\frac1{Nd}\sum_{i,o}|r_{io}|,\quad
 \sigma_i=\operatorname{sign}(r_i),\quad B_i=I-A_i.
\]
The negative student-gradient fields are
\begin{align*}
 F_{\rm input}&=\frac1d\E[J^\top B^\top\sigma],\\
 F_{\rm stop}&=\frac1d\E[J^\top\sigma],\\
 F_c&=\frac1d\E[J^\top\operatorname{sign}(r_c)].
\end{align*}
The $1/d$ factor implements the mean over all entries; here $d=2$.
We specify $\operatorname{sign}(0)=0$, without claiming a classical
derivative at zero. No smoothing or gradient rescaling was used.

At matched initialization with $c=q_{\theta_0}$, suppose all residual
components are nonzero. Continuity on a finite batch gives a neighborhood
where original and cached residuals share one fixed sign array $\sigma_0$.
For smooth model maps their stopped/cached fields agree throughout this
neighborhood, and the full derivative blocks satisfy
\[
 S_{\rm stop}=S_c=\frac1{Nd}\sum_{i,o}\sigma_{0,io}
 \nabla_\theta^2 q_{io},\qquad T_{\rm stop}=T_c=0.
\]
Thus the squared-loss difference $\E[J^\top A J]$ does not transfer to
this absolute-loss comparison. The weighted output Hessian remains;
it is the derivative of the residual sign that vanishes locally.
Matched deterministic trajectories remain equal while their residual
signs agree at shared states, but sign changes can invalidate this
agreement. No absolute-loss trajectories were run in this audit.

The original full input-gradient rule still differentiates $J$ and $A$
inside a sign region. In particular, for a reference perturbation $v$,
\[
 T_{\rm input}[v]=-\frac1d\E[J^\top(D_\phi A[v])^\top\sigma_0],
\]
which need not vanish. It would be incorrect to infer absence of local
reference feedback for all absolute-loss rules.

All five checkpoints pass the declared local checks without changing
finite-difference step $10^{-5}$. Every initial residual is nonzero;
the smallest magnitude is $3.01\times10^{-5}$. All 10,260 loss/field
derivative probes preserve signs. Independent product-rule assembly
matches gradients and full derivatives within $1.01\times10^{-11}$.
Stopped/cached blocks match exactly numerically, while full input-gradient
reference-block Frobenius norms range from $.449$ to $.847$ across all
five seeds. A separate closed-form tanh Hessian check agrees within
$8.33\times10^{-17}$. The coordinate probes do not certify a whole
neighborhood radius. This verifies a local calculus distinction in the
synthetic checkpoints, not a new stability theorem or an official benchmark.

\paragraph{Completed absolute-loss trajectories: an unresolved risk result.}
The unchanged protocol reused all five checkpoints, the 64/192 split,
and the three backward/input conditions with fixed reference or
$\tau=1$ EMA. All 60 trajectories and 90,000 updates completed at
$\eta=.02/.01$ to $T=20$, without smoothing, rescaling or retuning.
Labels remained evaluator-only, and every arm retained the same
128-number cache. The numerical gate required at most $5\%$ terminal
NMSE change under step halving in every primary stopped/cached condition,
measured as $|R_{\eta}-R_{\eta/2}|/\max(|R_{\eta/2}|,.01)$.
Fifteen of sixteen fail; the largest discrepancy is $152.28\%$.
The cached/EMA discrepancies alone range from $55.88\%$ to $93.14\%$.
Thus the primary risk inference is \emph{numerically unresolved},
not an established positive replication or negative effect finding.

For the EMA arms, the recorded terminal NMSE is:
\begin{center}
\begin{tabular}{rrrrr}
Seed & Original $.02$ & Cached $.02$ & Original $.01$ & Cached $.01$\\ \hline
0 & $22.27344$ & $.73696$ & $50.82051$ & $3.54438$\\
2 & $161.21032$ & $3.14709$ & $160.90436$ & $7.13295$\\
3 & $11.14135$ & $1.69339$ & $9.03343$ & $4.14852$\\
4 & $.83037$ & $.16184$ & $6.06859$ & $2.35862$
\end{tabular}
\end{center}
The previous $75\%$ excess-removal inequalities hold in all four seeds
at $.02$, but only seeds 0 and 2 at both steps. This is descriptive
because the numerical gate failed. Every cached/EMA endpoint remains
worse than initial accuracy. Six of eight full input-gradient secondary
checks also fail; those arms do not rescue the primary comparison.
All sensitivity-seed results remain available without entering its count.

The onset result is narrower and independently checked. In all 20
stopped/cached pairs, states remain exactly equal before the first
residual-sign disagreement. The weighted field difference is nonzero
at that event, and parameter states separate on the next update.
For a shared state at step $n_*$,
\[
 \theta_{{\rm stop},n_*+1}-\theta_{c,n_*+1}
 =\frac{\eta}{Nd}\sum_i J_i^\top
       (\sigma_{{\rm stop},i}-\sigma_{c,i}).
\]
This is an exact discrete identity; a general sign disagreement can
cancel under the weighted sum. Primary first sign disagreements occur
at times $.02$--$.04$ for the main step and $.01$--$.05$ for the fine
step. These are grid events, not converged continuous crossing times.
No residual is exactly zero at a sampled update, which does not exclude
crossings between updates. Sign changes alone do not explain terminal
risk or prove a particular nonsmooth limiting evolution.

An independent saved-data audit reproduces all 3,060 metric records and
targets exactly, checks 20 archives and 60 event snapshots, and verifies
the onset fields within $3.89\times10^{-16}$ and first-step separation
within $1.61\times10^{-16}$. Fixed-weight cached targets stay constant.
These checks validate the recorded algorithms, not step-size robustness.

\paragraph{Completed scalar boundary analysis.}
The solution concept matters for absolute discrepancy. As a declared
diagnostic, take $f_\theta(z)=h+\theta+az$, with $h=1$, $a=1/4$ and
$\theta_0=\phi_0=0$. Correct initial prediction $h$ is planted; this
is not an idempotence-trained model. Write $c=ah=1/4$ and
$e=\theta-\phi$. With cached input $q_0=h$, the residual is $c-e$.
The student-first discrete dynamics give exactly
\[
 e_{n+1}=\alpha\{e_n+\eta\operatorname{sign}(c-e_n)\},\qquad
 \phi_{n+1}-\phi_n=\frac{\eta}{\tau}e_{n+1},\quad
 \alpha=\frac{\tau}{\tau+\eta}.
\]
For a fixed reference, constant-step updates reach the residual root
only if it lies on their update grid; otherwise they enter a two-cycle.
At $\eta=1/10$, the cached example alternates between $1/5$ and $3/10$.

For moving reference and $0<c<\tau$, the formal pointwise lag ODE
$\dot e=\operatorname{sign}(c-e)-e/\tau$, with
$\operatorname{sign}(0)=0$, reaches $e=c$ at
$t_*=-\tau\log(1-c/\tau)$ but admits no absolutely continuous
continuation satisfying that ODE almost everywhere. Indeed, starting
at the boundary, $V=(e-c)^2/2$ satisfies
\[
 \dot V\le -(1-c/\tau)|e-c|-(e-c)^2/\tau\le0.
\]
It must remain zero, contradicting the pointwise value $\dot e=-c/\tau$.
This is an obstruction to that solution concept, not to the discrete
algorithm. Standard discontinuous-system theory distinguishes it from
the Filippov extension~\cite{cortes}.

With $\operatorname{Sign}(0)=[-1,1]$, the explicitly declared inclusion
$\dot e\in\operatorname{Sign}(c-e)-e/\tau$ has the unique solution
\[
 e(t)=\min\{\tau(1-e^{-t/\tau}),c\}.
\]
Monotonicity gives contraction between any two lag solutions; the
potential $|e-c|+e^2/(2\tau)$ places this scalar example within the
established maximal-monotone framework~\cite{benilanbrezis}.
After contact, the selected student velocity is $c/\tau$, so both
parameter copies keep moving while their discrepancy is zero.
Lag contraction therefore does not preserve the initially correct output.

An elementary bound now connects the exact cached discrete algorithm to
this inclusion. From the matched zero initialization, for every $\eta>0$
and grid time $t_n=n\eta$, with $C=1+c/\tau$,
\begin{align*}
 |e_n-e(t_n)|&\le C\eta,\\
 |\phi_n-\phi(t_n)|&\le\eta(Ct_n/\tau+c/\tau),\\
 |\theta_n-\theta(t_n)|&\le\eta\{C(1+t_n/\tau)+c/\tau\}.
\end{align*}
The proof bounds the pre-contact implicit-Euler lag error by $\eta/2$,
then uses the invariant interval
$[\alpha(c-\eta),\alpha(c+\eta)]$ and the exact reference-sum identity.
The parameter bounds are for fixed horizons, not uniform in time;
they do not certify derivatives, crossing times or neural accuracy.
Exact rational checks of 72 scalar updates at the sole prescribed
$\eta=1/10$ support the algebra. Full proofs and source locations are
saved in the research notes; no neural experiment was rerun.

\paragraph{Completed multivariate closure audit.}
The vector extension must preserve joint sign constraints. Let the actual
state be $z\in\R^p$, let $r:U\to\R^m$ be $C^1$, and let $M,b$ be
continuous, with $F=b+M\operatorname{sign}(r)$ and
$\operatorname{sign}(0)=0$. Here $m=Nd$; for the stopped student the
student block of $M$ is $J_q^\top/m$. Frozen reference weights are
parameters; moving reference weights belong to $z$, with their continuous
EMA drift in $b$. At $z_0$, let $\mathcal S(z_0)$ contain those joint
sign arrays in $\{-1,0,1\}^m$ that occur on a set of positive
$p$-dimensional measure in \emph{every} neighborhood of $z_0$. Then
the exact Filippov closure is
\[
 \mathcal K[F](z_0)=b_0+M_0\operatorname{co}\mathcal S(z_0).
\]
There are finitely many patterns: the nonessential ones can be removed
by one local null set, while every essential one persists after any such
removal. Continuity of $M,b$ proves the displayed identity. Zero entries
must be retained if a residual vanishes on positive measure.

Let $\mathcal B(r_0)$ fix the nonzero signs and replace each zero sign
independently by $[-1,1]$. Always
$\mathcal K[F](z_0)\subseteq b_0+M_0\mathcal B(r_0)$.
If $Z=\{j:r_j(z_0)=0\}$ and $D_zr_Z(z_0)$ has full row rank, local
submersion coordinates make every active sign orthant essential, so
equality holds. This sufficient condition is classical
Filippov calculus~\cite{padensastry,cortes}, not a new theorem.
Rank must be computed in the actual state space. It is not necessary:
the matrix $M_0$ may collapse the extra sign choices.

\paragraph{Why true partial gradients are different.}
For $C^1$ residuals, $D_1(z)=m^{-1}\sum_j|r_j(z)|$ is a finite maximum
of smooth signed sums. Its Clarke generalized gradient is exactly
\[
 \partial D_1(z_0)=m^{-1}D_zr(z_0)^\top\mathcal B(r_0).
\]
Let $P$ keep student coordinates and zero reference coordinates
($P=I$ for a frozen reference). The specified pointwise field agrees
almost everywhere with $b-P\nabla D_1$. The standard gradient-closure,
continuous-matrix and finite-maximum rules therefore give
\[
 \mathcal K[b-P\nabla D_1](z_0)=b_0-P\partial D_1(z_0).
\]
Thus the independent-sign \emph{velocity image} is exact for the original
full input-gradient and cached-input rules, without a rank assumption,
for smooth finite-batch maps and continuous reference drift.
For these rules $D_\theta r=(A-I)J_q$ and
$D_\theta r_c=-J_q$, respectively. The stopped field has a different
coefficient and does not inherit this conclusion. These are applications
of established calculus~\cite{padensastry,cortes}; they do not address
ReLU activation boundaries, Adam, solution uniqueness or useful adaptation.

\paragraph{An idempotent reference with a spurious interval-model path.}
Consider one input, two outputs, and the shared architecture
\[
 f_w(z)=h+w+Az,\qquad
 h=\begin{pmatrix}1\\-1\end{pmatrix},\quad
 A=\begin{pmatrix}0&0\\1&1\end{pmatrix}.
\]
Take $\phi=0$ fixed, $\theta=(u,v)$, $q_\theta=f_\theta(0)$ and
matched initialization $\theta_0=\phi_0=0$. Since $A^2=A$ and $Ah=0$,
$f_0\circ f_0=f_0$ globally. The evaluation target $h$ is planted and
never used in an update. The original residual and stopped field are
\[
 r=(-u,u)^\top,\qquad D_1=|u|,\qquad
 F_{\rm stop}=\tfrac12(-\operatorname{sign}u,\operatorname{sign}u)^\top.
\]
At $u=0$, its true closure is the segment
$\{(-\lambda/2,\lambda/2):|\lambda|\le1\}$, whereas independent signs
give the square $[-1/2,1/2]^2$. Every true solution obeys
$(u+v)'=0$ and $(u^2/2)'=-|u|/2$ almost everywhere. From the origin
it is uniquely stationary; the sign-zero discrete algorithm also remains
there for every step size. The enlarged square instead permits
$\theta(t)=(0,t/2)$, with zero discrepancy and evaluator risk
$R(t)=\norm{q_\theta-h}^2/2=t^2/8$. This apparent deterioration is an
artifact of the enlarged model, not a failure of the actual algorithm.

On the same model the full input-gradient closure is
$[-1,1]\times\{0\}$, exactly the independent-sign image despite rank
deficiency. Cached input gives $r_c=(-u,-v)$ and a full-rank residual
map. If the reference instead moves with state $\phi=(p,q)$, then
$r=(p-u,q+u)$ has rank two with respect to $(u,v,p,q)$, so the stopped
interval image is exact in that full state space. The frozen-reference
counterexample is not an EMA counterexample. Nor is it a bug in the
audited regression implementation, which retains the input gradient.

Exact-fraction checks verify the idempotence identities, all three fields
at 25 rational substitutions, the closure images, the invariant and
the analytic spurious path. No neural trajectories were rerun. A Lyapunov
inequality valid for \emph{every} velocity in an outer model remains
sufficient for the true closure, but an individual outer-model trajectory
need not be realizable. Neither closure equality nor this counterexample
resolves the failed neural numerical gate.

\paragraph{Completed paper decision and a narrower target.}
The nonsmooth closure results remain thesis background; they do not yet
support a first paper. Bolte--Pauwels already analyze automatic
differentiation using conservative derivative fields, including ordinary
nonsmooth networks~\cite{boltepauwels}. Their assumptions do not justify
an arbitrary stopped-output rule as a derivative of the same forward loss,
nor do they provide a guarantee for our constant-step EMA dynamics or
the audited Adam implementation.
Patsenker--Li--Kluger's BELAY already couples model and averaged weights
through direct parameter mixing~\cite{belay}. Its printed Algorithm~1
updates the reference from the old student; our reference uses the updated
student. General feedback or a damped-spring interpretation is consequently
insufficient novelty. Our prediction-input intervention is a different
operation, whose practical and theoretical value still needs validation.

The tested conditional question was: \emph{can computable finite-step
error bounds predict the output effect of prediction-input feedback and
certify the range over which that prediction is reliable?} The test uses the
existing smooth squared-loss comparison, whose early forecast succeeds
empirically but lacks a certified validity range. The intended result was
an explicit domain and output-error enclosure, with no evaluator labels
used to obtain it. Standard Taylor estimates are tools, not the novelty.

\paragraph{Completed certificate feasibility test.}
For a smooth joint update $\Phi$, let $c=\Phi(z_0)-z_0$,
$A=D\Phi(z_0)$ and $u_{n+1}=c+Au_n$, $u_0=0$.
If $D\Phi$ is $L$-Lipschitz on a ball of radius $R$, Taylor's
integral remainder gives the state-error recurrence
\[
 E_{n+1}=\norm{A}E_n+\tfrac L2(\norm{u_n}+E_n)^2,
 \qquad E_0=0.
\]
The induction requires $\norm{u_n}+E_n\le R$ at every expansion point.
For the RMS-scaled output map $Q$, a Hessian bound $H$ then gives
\[
 B_n=\norm{DQ(z_0)}E_n+\tfrac H2(\norm{u_n}+E_n)^2.
\]
Summing both arm bounds controls the output-difference forecast. These
are reproduced standard estimates, not new theory or a verified domain
for the trained networks. The saved update matrices use finite differences;
a positive certificate must additionally bound their coefficient errors
and numerical propagation defects. Those terms cannot be silently set to zero.

The fixed test used the five existing checkpoints, with seeds 0, 2, 3, 4
primary, both steps $.02/.01$, $T=2$, and the unchanged cache/input budgets.
The radius was twice the largest affine displacement through that horizon.
No neural adaptation, training or evaluator-label use was added.
The target was a valid bound at most $20\%$ of the predicted separation
in at least three primary seeds at both steps.

One explicit analytic construction, based on global tanh derivative and
weight-norm bounds, fails all eight primary conditions. Its output term
alone, even with zero state error, is $27.6\%$--$120.3\%$ of the predicted
separation. Its state-remainder term also exceeds the budget. Exact rational
comparisons and outward-rounded constants verify this failure; the total
calculation took under two CPU seconds. A separate necessary lower bound
on the best possible output-Hessian constant was inconclusive. Consequently
this result rejects the tested conservative construction, not all possible
certificates. No positive state-containment or network-error certificate
was obtained. The earlier empirical forecast accuracy remains unchanged.

The finite-step certificate route is therefore demoted in its present form.
Do not change the radius, norm, horizon, seeds or threshold to rescue it.
Paired or directional bounds remain untested, rather than becoming an
automatic replacement claim. Output separation still does not certify
current-task accuracy; the absolute-loss numerical failure also remains.

\paragraph{Completed forward/backward graph audit.}
Changing the reference's input value and changing its derivative are two
separate interventions. Let $z_v$ be either current $q_\theta$ or cached
$c=q_{\theta_0}$, let $b\in\{0,1\}$ specify its derivative with respect
to $q_\theta$, and write $A_v=D_zT_\phi(x,z_v)$ and
$r_v=T_\phi(x,z_v)-q_\theta$. The negative student directions are
\begin{align*}
 U^{(2)}_{v,b}&=\frac1N\sum_i J_i^\top(I-bA_{v,i})^\top r_{v,i},\\
 U^{(1)}_{v,b}&=\frac1{Nd}\sum_i
 J_i^\top(I-bA_{v,i})^\top\operatorname{sign}(r_{v,i}).
\end{align*}
These use the previous squared normalization and the official mean absolute
normalization, respectively. Current input with $b=1$ is the audited path;
current input with $b=0$ is the stopped comparison. Cached input normally
has $b=0$. The diagnostic $c+(q-\operatorname{sg}(q))$ instead has cached
forward value and $b=1$; it is a specified straight-through rule, not an
ordinary gradient of its forward loss or the official method.

At matched initialization, current and cached variants have equal gradients
within each fixed-$b$ pair. Equal first Adam updates additionally require
matching moment buffers, step counters and hyperparameters. The pinned
cell resets student weights per batch but carries optimizer state, so
matching weights alone is insufficient for comparisons across batches.

For absolute loss at a shared state, equal residual signs and equal $A_v$
make the connected current/cached directions equal. Identical reference
ReLU activation masks are sufficient for the latter equality. Thus their
first gradient difference requires a residual-sign or input-Jacobian
difference; either can still cancel after projection and averaging.
This discrete identity does not imply equal risk or a continuous crossing
time. It does not extend to squared loss merely from equal signs and masks.
Exact rational checks verify 96 local graph derivatives, 24 matched-gradient
and optimizer-moment identities, six vector decompositions and three scalar
witnesses. These are algebra checks, not a trained-network reproduction.

\paragraph{When does the second update change?}
The inspected tabular call takes exactly two Adam updates at learning rate
$10^{-4}$, with a fixed reference and no EMA. This permits a focused question:
does caching the forward input change the second update when the connected
backward path and initial optimizer state are preserved? The first update
is shared. At the common second evaluation, writing $s_v=\operatorname{sign}r_v$,
the connected-field difference is exactly
\[
 \Delta U=\frac1{Nd}\sum_iJ_i^\top\left[
 (I-A_{q,i})^\top(s_{q,i}-s_{c,i})
 +(A_{c,i}-A_{q,i})^\top s_{c,i}\right].
\]
This separates two contributions to the gradient difference, not to Adam's
nonlinear update or accuracy. The registered test includes all four graphs
and no adaptation, matched optimizer snapshots and $Nd$ cache scalars per
arm, with a cache created after each weight reset and discarded after the
batch. The source uses a ReLU model with default auxiliary input 50;
our tanh checkpoints cannot be silently substituted. These standard graph
identities sharpen the causal question but do not establish a novel
first-paper claim.

\paragraph{Feasibility result and artifact requirement.}
An existing CPU PyTorch runtime passes a data-free conformance check on
the source-sized 1291-parameter ReLU architecture. All four autograd
directions match their formulas; paired first updates agree when Adam
history is cloned. A negative control changes only that history and
changes the update. Weight resets preserve the moments and step counter.
These checks took 8.20 CPU seconds; they are not trained-model evidence.

The bounded local inventory found the existing 114-parameter tanh models
but no compatible tabular checkpoint or Boston Housing input file in the
proposal, adjacent research repository or chat-project folders. This is
not a machine-wide absence claim. The notebook normalizes features and
targets before splitting rows, so a checkpoint also needs its split and
scaling provenance. A new training-only normalization would be a declared
change, not an exact reproduction of that preprocessing.

The saved conditional plan selects batch size two, the first 20 complete
batches, four adaptation arms with two steps each, and no adaptation as a
control: 160 adaptation updates under a 60-CPU-second cap. It fixes numerical
tolerances and a measurable-effect criterion before any trained outputs
are viewed. The complete pinned repository tree and its tabular README
link provide a training notebook but no tabular checkpoint. The directly
linked dataset was saved at a fixed revision and content hash before
training values were inspected. This resolves data provenance for a new local
reproduction; it does not recover the authors' trained model.

\paragraph{Completed registered checkpoint preparation.}
The unchanged seed-zero run completed 30 epochs and 450 updates in
2.875 CPU seconds, below the 120-second cap. It used only the first 455
rows, training-only feature/target statistics, the specified ReLU/k50
architecture, two squared pretraining losses and batch size 32. Only
the final checkpoint was considered. Its training-only qualification is
\begin{align*}
 \mathrm{NMSE}_{\rm default}&=0.1919\le .75,\\
 \mathrm{NMSE}_{\rm label\ input}&=0.05521\le .25,\\
 \frac{\operatorname{std}(q)}{\operatorname{std}(y)}&=0.8699\ge .10.
\end{align*}
All three preregistered gates pass. These qualify a fitted, nonconstant
checkpoint for the diagnostic; they do not establish generalization or
useful adaptation. Training self-consistency MSE is nonzero ($0.01475$),
so no exact idempotence is claimed. An independent NumPy forward evaluation
of the saved tensors reproduces all reported metrics within $3.34\times10^{-16}$.
Saved scaling, split indices, tensor shapes, hashes and scheduler trace
were also checked without retraining.

Training-only normalization and CPU double precision are departures from the
notebook. The checkpoint is a local reproduction, not the authors' weights.

\paragraph{Two updates on held-out features.}
We compared a fresh reference input with the cached initial prediction while
keeping the gradient path through that input connected in both arms. The
cached arm uses the straight-through rule defined above. Both arms began
each batch with identical student weights and Adam history. We used the first
40 held-out features in order, without corruption or labels, in 20 pairs.
All four arms took two steps per pair (160 updates); only the current-connected
arm supplied optimizer history to the next pair. The no-adaptation baseline
retained the original predictions.

Five of 20 pairs met both thresholds, exactly the registered minimum:
a gradient difference exceeding $10^{-4}\max(\|g_q\|,\|g_c\|,10^{-8})$
and a final prediction RMS difference exceeding
$10^{-6}\max(\operatorname{RMS}(q_0),1)$. All five had a sign change in
at least one component of the absolute-loss residual. In four, the reference
input derivative was unchanged; in the fifth, one ReLU switched and the
derivative contribution partly offset the sign contribution. Prediction
differences in the five pairs ranged from $0.000278$ to $0.00948$ RMS
in normalized target units. In the other 15, the second gradients and
final predictions were identical. The largest componentwise error in the
gradient decomposition was $6.67\times10^{-16}$ or less.

\paragraph{The registered accuracy check fails.}
We then scored the locked post-update predictions on all 40 examples
using their held-out targets exclusively for evaluation. The protocol
required current-connected predictions to have lower mean squared error
than both the initial predictions and the cached-connected predictions.
Target scaling used the training statistics. The resulting normalized
mean squared errors were
\begin{center}
\begin{tabular}{lr}
\hline
Prediction rule & Mean squared error\\
\hline
No adaptation & 0.072213\\
Current input, connected derivative & 0.076643\\
Cached input, connected derivative & 0.076341\\
Current input, stopped derivative & 0.076390\\
Cached input, stopped derivative & 0.074749\\
\hline
\end{tabular}
\end{center}
Current-connected error rose by $6.14\%$ against no adaptation and
$0.40\%$ against cached-connected inputs. All four arms increased error.
The paired current-minus-cached error was exactly zero on the 15 batches
with identical predictions. Independent scalar summation verified the
saved errors and failed benefit criterion.

These results concern post-update predictions on the same 40 inputs
used for adaptation, one trained checkpoint and one ordered split.
The reference was fixed, student weights reset per batch, and canonical
Adam history carried across batches as specified above. The evidence
supports a mechanism for gradient separation alongside a measured accuracy
cost in this setting. Extension to changing distributions and retained
knowledge requires separately justified models and evaluation.
The current empirical branch closes with its failed accuracy criterion.
The two-source comparison and the resulting paper decision appear in
the thesis-scope section. The remaining theoretical question concerns
the information needed to certify the direction of the risk change.

A compatible $A\to B\to A$ stream remains conditional on a useful,
implementation-relevant mechanism and reliable numerical evaluation.

For that subsequent stream study, compare no adaptation, episodic reset,
carried student with a fixed
reference, carried student with EMA, and the explicit output-detachment
comparison. Select the closest established restoration/anchoring comparator
after the mechanism is identified. Report trainable parameters, reference
copies, source statistics, replay memory, teacher batch statistics, and
update counts; match these budgets or disclose differences. Report current
pre-update risk, post-update risk, discrepancy, reference lag, and boundary
prediction errors. On compatible recurring domains, also report earlier
risk and recovery relative to both no adaptation and a fresh source reset
under equal new-data/update budgets. Compare retention at matched current
risk so that slower learning is not mistaken for protection.

\paragraph{Four milestones, subject to the rejection gates.}
The constructed feedback family, useful fixed point, and discrete
certificates complete the first milestone. The first neural geometry gate
is also complete, with a negative result for strong rotation. The finite-time
gate now supports EMA amplification but rejects uniform backward-rule
ordering. The completed local decomposition rejects mismatch as a
necessary explanation, while the $L=J$ rates are not uniformly
accurate. The closest-prior-art claim map is complete and rules out
generic feedback failure as a sufficient paper claim. The predeclared
prediction-input control is complete and passes in all primary seeds,
while failing to preserve initial accuracy. The early output forecast
now passes in all primary seeds, with a documented long-horizon failure.
The local absolute-loss audit is also complete: it separates the
stopped/cached sign-region identity from nonzero full-gradient reference
feedback. The finite-time absolute-loss protocol is complete, with a
verified separation onset but a failed numerical gate for risk. The
scalar nonsmooth analysis supplies both an obstruction to the pointwise
ODE and a bound for the correct inclusion. The multivariate closure audit
now gives an exact joint-sign formula, a standard rank condition, and
a checked idempotent-reference counterexample for the stopped rule.
Its smooth partial-gradient corollary does not cover the official
ReLU/Adam graph; the failed neural protocol has not been retuned.
The final milestone is reproduction of one selected implementation graph
and robustness to its loss and optimizer. An affine approximation alone is not a novel theorem.
The bounded claim decision is complete: the closure branch remains
background, and the concrete finite-step certificate construction fails
its nonvacuity gate. The failure is retained without retuning. The completed
graph audit supplies a discrete equality criterion and a two-update causal
question. CPU graph conformance and data provenance are resolved for a
new local reproduction. The registered final checkpoint passes its
training-only qualification. The two-update diagnostic meets its separation
criterion. The registered same-batch accuracy check fails, with all four
adaptation variants increasing mean squared error over the baseline.
No publication outcome is assumed.

\paragraph{Existing work and limits.}
The earlier scalar learned-predictor certificate remains a checked local
research artifact, useful for proof practice. Its PEIRA-inspired predictor
tracks an optimizer, whereas the present reference tracks the student;
the earlier result does not prove anything about EMA adaptation. Likewise,
function-space stop-gradient convergence results require their own operator
and regularity assumptions~\cite{stopgrads}. The project currently has a source/code audit, the appendix calculation,
constructed local stability results, and reproducible synthetic neural
geometry, finite-time response, local feedback decomposition, and a
paired frozen-input intervention with a checked early output forecast,
plus local and finite-time absolute-loss audits, a scalar nonsmooth
baseline, and a faithful-closure audit that delimit its transfer.
It has no general trained-network stability theorem,
distribution-drift guarantee, official-implementation reproduction,
or verified novelty claim.

\medskip
\noindent{\small\textit{Acknowledgments.} This proposal has developed through ongoing discussion and collaboration with Shoham Grunblat.}

\appendix
\section{A checked affine diagnostic, not a paper claim}
Take unit features $x_A,x_B\in\R^2$ with $0<|c|<1$, where
$c=x_A^\top x_B$. Let $f_\theta(x,z)=\theta^\top x+a(x)z$,
$a(x_A)=0$, and $a(x_B)=1-b$ for $b>1$. Start with
$\theta_0=\phi_0$, let $y_A=\theta_0^\top x_A$ and $y_B=0$, and use
these labels only for evaluation. During the $B$ phase, write
$u=\theta^\top x_B$, $v=\phi^\top x_B$, and $u_0\ne0$. Then
\begin{equation}
 \dot u=k(v-bu),\qquad \tau\dot v=u-v,
 \qquad k=b\ \text{for }G_{\rm input},\quad k=1\ \text{for }G_{\rm stop}.
 \label{eq:toy}
\end{equation}
The two reduced eigenvalues are negative for every $\tau,k>0$; the
orthogonal student component is conserved. In fact $u$ is a positive
weighted sum of two decaying exponentials times $u_0$. Consequently
\begin{equation}
 D=\tfrac12(v-bu)^2\downarrow0,\qquad
 R_B=\tfrac12u^2\downarrow0,\qquad
 R_A=\tfrac12c^2(u_0-u)^2\uparrow\tfrac12c^2u_0^2.
 \label{eq:toyresult}
\end{equation}
Both tasks are jointly realizable. The limiting state is
$\theta_\infty=\phi_\infty=\theta_0-u_0x_B$, independent of the finite
EMA timescale. Returning to $A$ from that limit provides zero consistency
signal and hence no recovery. A finite-time return can recover partially
through reference lag. Full derivations and verification are stored in
the research notes; 40 parameter combinations at 502 times each passed
the closed-form checks, with maximum ODE residual about $1.1\times10^{-13}$.
Numerical checks support the algebra; they are not a neural experiment.

This diagnostic shows that useful current adaptation can coexist with loss
on a compatible earlier task, and that a slow EMA need not preserve it
indefinitely. Its two backward rules only change speed and have identical
retention at equal $B$ progress. The deliberately chosen feedback and label
are not established empirical IT$^3$ properties. The result motivates the
vector study and its falsification gates; it is not an independent novelty
claim. Protecting the old feature direction would require retained
information and overlaps with established projected-update ideas~\cite{ogd}.

\begin{thebibliography}{99}\small
\bibitem{it3} N. Durasov et al. \emph{IT$^3$: Idempotent Test-Time Training}.
ICML, PMLR 267:14867--14883, 2025.
\href{https://proceedings.mlr.press/v267/durasov25a.html}{Published record};
\href{https://arxiv.org/abs/2410.04201}{arXiv:2410.04201, v1/v2}.
\bibitem{ttt} Y. Sun et al. \emph{Test-Time Training with Self-Supervision
for Generalization under Distribution Shifts}. ICML, 2020.
\href{https://proceedings.mlr.press/v119/sun20b.html}{PMLR 119}.
\bibitem{cotta} Q. Wang et al. \emph{Continual Test-Time Domain Adaptation}.
CVPR, 2022.
\href{https://openaccess.thecvf.com/content/CVPR2022/html/Wang_Continual_Test-Time_Domain_Adaptation_CVPR_2022_paper.html}{Primary paper}.
\bibitem{eata} S. Niu et al. \emph{Efficient Test-Time Model Adaptation
without Forgetting}. ICML, 2022.
\href{https://proceedings.mlr.press/v162/niu22a.html}{PMLR 162}.
\bibitem{rmt} M. Doebler, R. A. Marsden, and B. Yang. \emph{Robust Mean
Teacher for Continual and Gradual Test-Time Adaptation}. CVPR, 2023.
\href{https://arxiv.org/abs/2211.13081}{arXiv:2211.13081}.
\bibitem{petta} T.-H. Hoang, D. M. Vo, and M. N. Do. \emph{Persistent
Test-time Adaptation in Recurring Testing Scenarios}. NeurIPS, 2024.
\href{https://arxiv.org/abs/2311.18193}{arXiv:2311.18193}.
\bibitem{fixedteacher} D. Sojka et al. \emph{Rethinking the Teacher-Student
Framework for Test-Time Adaptation}. Preprint, 2026.
\href{https://arxiv.org/abs/2609.02507}{arXiv:2609.02507}.
\bibitem{desa} A. Belal et al. \emph{DESA-TTA: Dynamic EMA and Source
Anchoring for Test-Time Adaptation}. Preprint, 2026.
\href{https://arxiv.org/abs/2609.01795}{arXiv:2609.01795}.
\bibitem{realign} Z.-Y. Zhang et al. \emph{Test-time Adaptation in
Non-stationary Environments via Adaptive Representation Alignment}.
NeurIPS, 2024.
\href{https://papers.neurips.cc/paper_files/paper/2024/hash/abe31a12e83111fdf2cfd54deed5a2ce-Abstract-Conference.html}{Proceedings}.
\bibitem{recovery} Z. Zhou et al. \emph{On the Learnability of Test-Time
Adaptation: A Recovery Complexity Perspective}. ICML, 2026.
\href{https://arxiv.org/abs/2605.28057}{arXiv:2605.28057}.
\bibitem{ogd} M. Farajtabar et al. \emph{Orthogonal Gradient Descent for
Continual Learning}. AISTATS, 2020.
\href{https://proceedings.mlr.press/v108/farajtabar20a.html}{PMLR 108}.
\bibitem{spurious} V. Palit et al. \emph{When Forgetting is not Catastrophic:
On the Mechanics of Spurious Forgetting}. Preprint, 2026.
\href{https://arxiv.org/abs/2610.08718}{arXiv:2610.08718}.
\bibitem{stopgrads} M. W. Shen et al. \emph{How I learned to stop worrying
and love StopGrads: Stationarity, Convergence, and a case study on Flow Map
Learning}. Preprint, 2026.
\href{https://arxiv.org/abs/2609.16222}{arXiv:2609.16222}.
\bibitem{boltepauwels} J. Bolte and E. Pauwels.
\emph{Conservative set valued fields, automatic differentiation,
stochastic gradient methods and deep learning}.
\href{https://arxiv.org/abs/1909.10300v4}{Inspected arXiv v4, 2020},
Section 5.2, Theorem 8 and Corollary 5.
\bibitem{belay} J. Patsenker, H. Li and Y. Kluger.
\emph{Exponential weight averaging as damped harmonic motion}.
ICML Workshop on New Frontiers in Learning, Control, and Dynamical Systems, 2023.
\href{https://arxiv.org/abs/2310.13854v1}{arXiv v1}; Section 3 and Algorithm 1.
\bibitem{padensastry} B. E. Paden and S. S. Sastry.
\emph{A Calculus for Computing Filippov's Differential Inclusion with
Application to the Variable Structure Control of Robot Manipulators}.
25th IEEE Conference on Decision and Control, 1986.
Theorem 1, printed p.579; inspected conference version.
\href{https://doi.org/10.1109/CDC.1986.267369}{Conference DOI}.
\bibitem{cortes} J. Cort\'es. \emph{Discontinuous Dynamical Systems:
A tutorial on solutions, nonsmooth analysis, and stability}.
IEEE Control Systems Magazine 28(3):36--73, 2008.
\href{https://arxiv.org/abs/0901.3583v1}{Inspected arXiv v1};
\href{https://doi.org/10.1109/MCS.2008.919306}{Published DOI}.
\bibitem{benilanbrezis} P. B\'enilan and H. Br\'ezis.
\emph{Solutions faibles d'\'equations d'\'evolution dans les espaces de Hilbert}.
Annales de l'Institut Fourier 22(2):311--329, 1972.
\href{https://aif.centre-mersenne.org/item/AIF_1972__22_2_311_0/}{Primary article},
Definition 1.1, Lemma 1.1 and Proposition 1.1.
\end{thebibliography}
\end{document}
